\documentclass[11pt]{article}

\usepackage[a4paper,top=19mm,bottom=19mm,left=24mm,right=24mm]{geometry}
\usepackage[T1]{fontenc}
\usepackage{newtxtext}
\usepackage{microtype}
\usepackage[hidelinks]{hyperref}

\pagestyle{empty}
\setlength{\parindent}{0pt}
\setlength{\parskip}{4.5pt}

\begin{document}
\raggedright
\hyphenpenalty=10000
\exhyphenpenalty=10000

25 July 2026

Editorial Office\\
\textit{Knowledge-Based Systems}

\textbf{Re: Submission of ``Evidence contracts for certified knowledge
acquisition in retrieval-augmented generation systems''}

Dear Editor,

On behalf of all authors, we submit the above manuscript for consideration in
\textit{Knowledge-Based Systems}.

The manuscript addresses a knowledge-representation problem that arises when
retrieval-augmented generation systems move from evaluation to deployment:
how should evidence obligations be represented, and how can a system reason
soundly about whether an adaptive acquisition policy satisfies them? We
introduce \emph{evidence contracts}, a machine-checkable representation that
links answer quality, citation completeness, evidence freshness, and latency
to admissible violation rates and a certification confidence. ContractRAG
turns this representation into an auditable deployment decision.

The method represents heterogeneous acquisition plans as a rewrite lattice.
Risk vouchers encode empirical knowledge about the effects of plan rewrites
and compose without independence assumptions. Learned sufficiency models use
runtime evidence to choose when a query should acquire stronger knowledge,
while a distribution-free fixed-sequence test is the sole source of
certification. We prove finite-sample contract validity regardless of
candidate count or score quality, characterize cost consistency, and extend
the certificate through knowledge drift with an anytime-valid e-process and
repeated recertification.

The empirical study covers four heterogeneous knowledge-intensive benchmarks,
four large language model families, and two serving backends. Across 1,000
calibration draws, certified policies exceeded their contracts in at most
0.8\% of draws under a 10\% error budget, whereas uncertified tuning methods
reached 64\% in the same evaluation. The selected policies reduced model cost
by up to 20$\times$ relative to the strongest fixed pipeline, or
10--18$\times$ after audit costs were included. The released execution
matrices and code reproduce the reported results without further model calls.

This contribution fits the journal's emphasis on knowledge representation,
reasoning, and systems that support accountable decisions. Its principal
claim is not merely that a RAG pipeline becomes cheaper. It is that explicit
knowledge obligations can be represented and used to derive a statistically
sound deployment decision; the cost savings follow from that certified
reasoning process.

This manuscript is original, has not been published previously, and is not
under consideration by another journal. All authors have read and approved the
manuscript and agree with its submission to \textit{Knowledge-Based Systems}.
The authors declare no competing interests.

Thank you for considering our manuscript. We would be grateful for the
opportunity to have it evaluated for publication in
\textit{Knowledge-Based Systems}.

Sincerely,

\textbf{Weiwei Fu and Jinjiang Cui}\\
Co-corresponding authors, on behalf of all authors\\
Suzhou Institute of Biomedical Engineering and Technology\\
Chinese Academy of Sciences, Suzhou, China\\
Email: \href{mailto:fuww@sibet.ac.cn}{fuww@sibet.ac.cn};
\href{mailto:cuijj@sibet.ac.cn}{cuijj@sibet.ac.cn}

\end{document}
